\documentclass[12pt]{article}
\usepackage[a3paper, margin=1in]{geometry}
\usepackage{amsmath, amssymb, amsthm, ulem}

\begin{document}

\section*{Domácí úkol 3.1}

Vypracoval: Daniel \textit{"Randál"} Ransdorf \hfill Podpis: \rule{4cm}{0.4pt}

\begin{enumerate}
  \item \textbf{Řešení}:
    $\gamma$ označme jako $0$, protože $\forall x \in T : x \cdot \gamma = \gamma$.
    Poté označme $\beta$ jako $1$, protože je to jediný prvek, který
    splňuje $\forall x \in T \setminus \{\gamma\} : x \cdot \beta = x$.
    Dále z tabulky vyvoďme opačné a inverzní prvky:
    \begin{align*}
      -\alpha &= \alpha \\
      -\beta &= \beta \\
      -\gamma &= \gamma \\
      -\delta &= \delta \\
      \alpha^{-1} &= \delta \\
      \beta^{-1} &= \beta \\
      \delta^{-1} &= \alpha
    \end{align*}

    \[
      \left(\begin{array}{cccc|c}
      \alpha & \alpha & \delta & \alpha & \delta \\
      \beta & \delta & \beta & \delta & \delta \\
      \gamma & \beta & \alpha & \beta & \delta \\
      \end{array}\right)
      \overset{\text{dosazení 0,1}}{\;\sim\;}
      \left(\begin{array}{cccc|c}
      \alpha & \alpha & \delta & \alpha & \delta \\
      1 & \delta & 1 & \delta & \delta \\
      0 & 1 & \alpha & 1 & \delta \\
      \end{array}\right)
    \]
    \[
      \overset{\text{1.ř } (\cdot \delta)}{\;\sim\;}
      \left(\begin{array}{cccc|c}
      1 & 1 & \alpha & 1 & \alpha \\
      1 & \delta & 1 & \delta & \delta \\
      0 & 1 & \alpha & 1 & \delta \\
      \end{array}\right)
      \overset{\text{2.ř }(+\text{1.ř})}{\;\sim\;}
      \left(\begin{array}{cccc|c}
      1 & 1 & \alpha & 1 & \alpha \\
      0 & \alpha & \delta & \alpha & 1 \\
      0 & 1 & \alpha & 1 & \delta \\
      \end{array}\right)
      \overset{\text{2.ř } \cdot \delta}{\;\sim\;}
      \left(\begin{array}{cccc|c}
      1 & 1 & \alpha & 1 & \alpha \\
      0 & 1 & \alpha & 1 & \delta \\
      0 & 1 & \alpha & 1 & \delta \\
      \end{array}\right)
    \]
    \[
      \overset{\text{3.ř } (+\text{2.ř})}{\;\sim\;}
      \left(\begin{array}{cccc|c}
      1 & 1 & \alpha & 1 & \alpha \\
      0 & 1 & \alpha & 1 & \delta \\
      0 & 0 & 0 & 0 & 0 \\
      \end{array}\right)
      \overset{\text{1.ř } (+ \text{2.ř})}{\;\sim\;}
      \left(\begin{array}{cccc|c}
      1 & 0 & 0 & 0 & 1 \\
      0 & 1 & \alpha & 1 & \delta \\
      0 & 0 & 0 & 0 & 0
      \end{array}\right)
      \;\sim\;
      \left(\begin{array}{cccc|c}
      1 & 0 & 0 & 0 & 1 \\
      0 & 1 & \alpha & 1 & \delta \\
      0 & 0 & 0 & 0 & 0 \\ \hline
      0 & 0 & 1 & 0 & t_1 \\
      0 & 0 & 0 & 1 & t_2
      \end{array}\right)
      \overset{\text{4.ř } \cdot \alpha}{\;\sim\;}
      \left(\begin{array}{cccc|c}
      1 & 0 & 0 & 0 & 1 \\
      0 & 1 & \alpha & 1 & \delta \\
      0 & 0 & 0 & 0 & 0 \\ \hline
      0 & 0 & \alpha & 0 & \alpha t_1 \\
      0 & 0 & 0 & 1 & t_2
      \end{array}\right)
    \]
    \[
      \overset{\text{2.ř } (+\text{4.ř})}{\;\sim\;}
      \left(\begin{array}{cccc|c}
      1 & 0 & 0 & 0 & 1 \\
      0 & 1 & 0 & 1 & \delta + \alpha t_1 \\
      0 & 0 & 0 & 0 & 0 \\ \hline
      0 & 0 & \alpha & 0 & \alpha t_1 \\
      0 & 0 & 0 & 1 & t_2
      \end{array}\right)
      \overset{\text{2.ř } (+\text{5.ř})}{\;\sim\;}
      \left(\begin{array}{cccc|c}
      1 & 0 & 0 & 0 & 1 \\
      0 & 1 & 0 & 0 & \delta + \alpha t_1 + t_2 \\
      0 & 0 & 0 & 0 & 0 \\ \hline
      0 & 0 & \alpha & 0 & \alpha t_1 \\
      0 & 0 & 0 & 1 & t_2
      \end{array}\right)
      \overset{\text{4.ř } (\cdot \delta)}{\;\sim\;}
      \left(\begin{array}{cccc|c}
      1 & 0 & 0 & 0 & 1 \\
      0 & 1 & 0 & 0 & \delta + \alpha t_1 + t_2 \\
      0 & 0 & 0 & 0 & 0 \\ \hline
      0 & 0 & 1 & 0 & t_1 \\
      0 & 0 & 0 & 1 & t_2
      \end{array}\right)
    \]

    \[
    K = \left\{ 
        \left(\begin{array}{c} 1 \\ \delta + \alpha t_1 + t_2 \\ t_1 \\ t_2 \end{array}\right) 
        \quad \middle| \quad t_1,t_2 \in T
      \right\} = 
      \left\{ 
        \left(\begin{array}{c} 1 \\ \delta \\ 0 \\ 0 \end{array}\right) +
        t_1\left(\begin{array}{c} 0 \\ \alpha \\ 1 \\ 0 \end{array}\right) +
        t_2\left(\begin{array}{c} 0 \\ 1 \\ 0 \\ 1 \end{array}\right)
        \quad \middle| \quad t_1,t_2 \in T
      \right\}
    \]
    Zpětným nahrazením $0,1$ řeckými symboly dostaneme řešení:
    \[
    K = \underline{\underline{\left\{ 
        \left(\begin{array}{c} \beta \\ \delta \\ \gamma \\ \gamma \end{array}\right) +
        t_1\left(\begin{array}{c} \gamma \\ \alpha \\ \beta \\ \gamma \end{array}\right) +
        t_2\left(\begin{array}{c} \gamma \\ \beta \\ \gamma \\ \beta \end{array}\right)
        \quad \middle| \quad t_1,t_2 \in T
      \right\}}}
    \]

    Protože prostor řešení závisí na dvou volných parametrech, počet všech řešení bude velikost množiny kartézského součinu $T \times T$ (určuje počet dvojic $t_1,t_2$),
    čili $|T \times T| = |T|^2 = 4^2 = \underline{\underline{16}}$.

\end{enumerate}

\end{document}
